\documentclass[10pt,a4paper]{article}
\usepackage[utf8]{inputenc}
\usepackage[vietnamese]{babel}
\usepackage{amsmath,amssymb,amsfonts}
\usepackage{geometry}
\geometry{a4paper, left=10mm, right=10mm, top=10mm, bottom=10mm}
\usepackage{tikz}
\usetikzlibrary{shapes,arrows.meta,positioning,calc}
\usepackage[most]{tcolorbox}
\usepackage{array,booktabs,multicol}
\usepackage{fontawesome}
\usepackage{helvet}
\renewcommand{\familydefault}{\sfdefault}
\usepackage{parskip}

\definecolor{primary}{RGB}{112, 48, 160}
\definecolor{teal}{RGB}{0, 128, 128}
\definecolor{orange}{RGB}{230, 108, 44}
\definecolor{green}{RGB}{46, 139, 87}
\definecolor{rose}{RGB}{192, 0, 0}

\pagestyle{empty}

\begin{document}

% --- HEADER BANNER ---
\begin{tcolorbox}[
  enhanced,
  colback=violet!8!white,
  colframe=violet!80!black,
  arc=2.5mm, boxrule=1pt,
  left=3mm, right=3mm, top=2mm, bottom=2mm
]
  \begin{tabular*}{\textwidth}{@{\extracolsep{\fill}} lr}
    {\small\bfseries\color{violet!85!black}\faBook\ SỔ TAY KIẾN THỨC TOÁN 2} & 
    {\small\bfseries\color{teal!85!black}\faGraduationCap\ CHƯƠNG TRÌNH GDPT 2018}
  \end{tabular*}
  \vspace{-1mm}
  \begin{center}
    \rule{0.98\textwidth}{0.6pt}
  \end{center}
  \vspace{-2mm}
  \centering
  {\LARGE\bfseries\color{violet!90!black}\faBookmark\ BÀI 73: ÔN TẬP ĐO LƯỜNG}\\[1mm]
  {\itshape\footnotesize\color{black!75} Tổng hợp kiến thức độ dài, khối lượng, dung tích, thời gian và tiền tệ}
\end{tcolorbox}

\vspace{-2mm}

\noindent
\begin{minipage}[t]{0.49\textwidth}\vspace{0pt}
\begin{tcolorbox}[
  enhanced,
  colback=violet!4!white,
  colframe=violet!75!black,
  arc=2mm, boxrule=0.8pt,
  left=2mm, right=2mm, top=1.8mm, bottom=1.8mm,
  title={\footnotesize\bfseries\color{white}\faKey\ 1. Bảng đơn vị đo lường toàn cấp lớp 2}
]
{\footnotesize
\begin{itemize}\setlength{\itemsep}{1pt}\setlength{\leftskip}{-2.5mm}
  \item \textbf{Độ dài:}
  $$\mathbf{1\,dm = 10\,cm}; \quad \mathbf{1\,m = 10\,dm = 100\,cm}; \quad \mathbf{1\,km = 1\,000\,m}$$
  \item \textbf{Khối lượng:} Ki-lô-gam ($\mathbf{kg}$) - dùng cân đĩa, cân đồng hồ để đo.
  \item \textbf{Dung tích:} Lít ($\mathbf{l}$) - dùng ca $1\,l$ để đong nước, sữa.
  \item \textbf{Thời gian:}
  \begin{itemize}\setlength{\itemsep}{0.5pt}\setlength{\leftskip}{-2mm}
    \item $1\text{ ngày} = 24\text{ giờ}$; $1\text{ giờ} = 60\text{ phút}$.
    \item $1\text{ tuần lễ} = 7\text{ ngày}$ (Thứ Hai $\to$ Chủ nhật).
    \item Các tháng có $31$ ngày: 1, 3, 5, 7, 8, 10, 12.
    \item Các tháng có $30$ ngày: 4, 6, 9, 11.
    \item Tháng 2 có $28$ hoặc $29$ ngày.
  \end{itemize}
\end{itemize}
}
\end{tcolorbox}
\end{minipage}
\hfill
\begin{minipage}[t]{0.49\textwidth}\vspace{0pt}
\begin{tcolorbox}[
  enhanced,
  colback=teal!4!white,
  colframe=teal!75!black,
  arc=2mm, boxrule=0.8pt,
  left=2mm, right=2mm, top=1.8mm, bottom=1.8mm,
  title={\footnotesize\bfseries\color{white}\faSliders\ 2. Kĩ năng xem đồng hồ và lịch}
]
{\footnotesize
\begin{itemize}\setlength{\itemsep}{1pt}\setlength{\leftskip}{-2.5mm}
  \item \textbf{Đọc giờ đúng:} Kim dài chỉ số $12$, kim ngắn chỉ số nào đọc giờ số đó (ví dụ: $8$ giờ đúng).
  \item \textbf{Đọc giờ rưỡi ($30$ phút):} Kim dài chỉ số $6$, kim ngắn ở giữa hai số (ví dụ: $8$ giờ $30$ phút hay $8$ giờ rưỡi).
  \item \textbf{Đọc giờ có 15 phút:} Kim dài chỉ số $3$.
  \item \textbf{Xem tờ lịch tháng:} Biết ngày nào trong tuần, ngày đầu tháng, ngày cuối tháng, số ngày trong tháng.
\end{itemize}
}
\end{tcolorbox}
\end{minipage}

\vspace{-1.5mm}
\begin{tcolorbox}[
  enhanced,
  colback=orange!4!white,
  colframe=orange!80!black,
  arc=2mm, boxrule=0.8pt,
  left=2mm, right=2mm, top=1.5mm, bottom=1.5mm,
  title={\footnotesize\bfseries\color{white}\faSitemap\ 3. Sơ đồ các đại lượng thực tế trong đời sống}
]
\centering
\begin{tikzpicture}[x=1.65cm, y=0.75cm]
  \tikzset{arr/.style={-{Stealth[scale=0.8]}}, every node/.append style={align=center}}
  \node[rectangle, rounded corners=2pt, draw=violet!80!black, fill=violet!15, font=\scriptsize\bfseries] (d) at (0,0.6) {ĐỘ DÀI\\$cm, dm, m, km$};
  \node[rectangle, rounded corners=2pt, draw=teal!80!black, fill=teal!15, font=\scriptsize\bfseries] (k) at (0,-0.6) {SỨC CHỨA \& NẶNG\\$kg$ và $l$};
  \node[rectangle, rounded corners=2pt, draw=orange!80!black, fill=orange!15, font=\scriptsize\bfseries] (t) at (3.5,0.6) {THỜI GIAN\\Giờ, ngày, tháng};
  \node[rectangle, rounded corners=2pt, draw=green!70!black, fill=green!15, font=\scriptsize\bfseries] (m) at (3.5,-0.6) {TIỀN TỆ\\$100 \dots 1\,000$ đồng};

  \draw[arr, thick] (d) -- (t);
  \draw[arr, thick] (k) -- (m);
\end{tikzpicture}
\end{tcolorbox}

\vspace{-1.5mm}
\noindent
\begin{minipage}[t]{0.49\textwidth}\vspace{0pt}
\begin{tcolorbox}[
  enhanced,
  colback=green!4!white,
  colframe=green!60!black,
  arc=2mm, boxrule=0.8pt,
  left=2mm, right=2mm, top=1.8mm, bottom=1.8mm,
  title={\footnotesize\bfseries\color{white}\faCheckCircle\ 4. Bài tập thực hành mẫu}
]
{\footnotesize
\textbf{1. Đổi đơn vị và tính toán:}\\
- $4\,m\ 6\,dm = 40\,dm + 6\,dm = \mathbf{46\,dm}$.\\
- $45\,kg + 38\,kg = \mathbf{83\,kg}$.\\
- $10\,l : 2 = \mathbf{5\,l}$.\\
\textbf{2. Bài toán thời gian:}\\
Buổi học bắt đầu lúc $8$ giờ sáng và kết thúc lúc $10$ giờ sáng. Buổi học kéo dài trong:\\
$10 - 8 = \mathbf{2}$ (giờ).\\
\textbf{3. Xem lịch:} Nếu ngày $2$ tháng 9 là thứ Ba thì ngày $9$ tháng 9 là \textbf{thứ Ba} (cộng thêm $7$ ngày).
}
\end{tcolorbox}
\end{minipage}
\hfill
\begin{minipage}[t]{0.49\textwidth}\vspace{0pt}
\begin{tcolorbox}[
  enhanced,
  colback=rose!4!white,
  colframe=rose!70!black,
  arc=2mm, boxrule=0.8pt,
  left=2mm, right=2mm, top=1.8mm, bottom=1.8mm,
  title={\footnotesize\bfseries\color{white}\faExclamationTriangle\ 5. Cạm bẫy \& Mẹo xem lịch}
]
{\footnotesize
\begin{itemize}\setlength{\itemsep}{1.5pt}\setlength{\leftskip}{-2.5mm}
  \item \textbf{Quy tắc cộng 7 ngày:} Cứ cộng thêm hoặc trừ đi $7$ ngày thì thứ trong tuần không thay đổi.
  \item \textbf{Mẹo nắm bàn tay xem ngày tháng:} Các khớp lồi lên có $31$ ngày, các chỗ lõm xuống có $30$ ngày (riêng tháng 2 có $28$ hoặc $29$ ngày).
  \item Khi làm phép tính luôn viết đơn vị đo kèm theo.
\end{itemize}
}
\end{tcolorbox}
\end{minipage}


\end{document}
